\documentclass[12pt]{article}
\usepackage{amsmath,amssymb,amsfonts}
\usepackage[margin=1in]{geometry}

\begin{document}

\section*{Chapter 2, Problem 3}

\subsection*{Problem Statement}
\textit{Geodesics on a sphere.} On a sphere of fixed radius $a$, the line element is given by
\[
  d\ell^2 = a^2\,d\theta^2 + a^2\sin^2\theta\,d\phi^2.
\]
Determine the equation of the curve that extremizes the distance between two points. Is the extremum a minimum, a maximum, or what? [Hint: Let $\theta$ be the independent variable so the integrand of the integral to be extremized is independent of the dependent variable $\phi$.] Now prove: if $\theta$ and $\phi$ are constants, where $\phi = a\theta$ and $\theta = a\phi$, then change from the spherical coordinates to Cartesian coordinates to see that this is the equation of a plane through the origin; its intersection with the sphere $r = a$ is the geodesic---a great circle---on the surface of the sphere.

\subsection*{Solution}

The arc length on the sphere is
\[
  L = \int a\sqrt{1 + \sin^2\theta\,\phi'^2}\,d\theta,
\]
where $\phi' = d\phi/d\theta$. The integrand $f(\theta, \phi') = a\sqrt{1 + \sin^2\theta\,\phi'^2}$ does not depend on $\phi$ explicitly. Therefore, by the Euler--Lagrange equation:
\[
  \frac{\partial f}{\partial \phi} = 0
  \quad \Longrightarrow \quad
  \frac{d}{d\theta}\frac{\partial f}{\partial \phi'} = 0
  \quad \Longrightarrow \quad
  \frac{\partial f}{\partial \phi'} = \text{const}.
\]
Computing:
\[
  \frac{\partial f}{\partial \phi'} = \frac{a\sin^2\theta\,\phi'}{\sqrt{1 + \sin^2\theta\,\phi'^2}} = C,
\]
where $C$ is a constant. Squaring:
\[
  a^2\sin^4\theta\,\phi'^2 = C^2(1 + \sin^2\theta\,\phi'^2),
\]
\[
  \phi'^2(a^2\sin^4\theta - C^2\sin^2\theta) = C^2,
\]
\[
  \phi'^2 = \frac{C^2}{\sin^2\theta(a^2\sin^2\theta - C^2)}.
\]
Separating variables:
\[
  d\phi = \frac{C\,d\theta}{\sin\theta\sqrt{a^2\sin^2\theta - C^2}}.
\]
Let $k = C/a$. Then:
\[
  d\phi = \frac{k\,d\theta}{\sin\theta\sqrt{\sin^2\theta - k^2}}.
\]
Substituting $u = \cot\theta$, so $du = -d\theta/\sin^2\theta$ and $\sin^2\theta = 1/(1+u^2)$:
\[
  d\phi = \frac{-k\,du}{\sqrt{1 - k^2(1+u^2)}}
  = \frac{-k\,du}{\sqrt{(1-k^2) - k^2 u^2}}.
\]
Let $\beta^2 = (1-k^2)/k^2$. Then:
\[
  d\phi = \frac{-du}{\sqrt{\beta^2 - u^2}}.
\]
Integrating:
\[
  \phi - \phi_0 = \arccos\!\left(\frac{u}{\beta}\right)
  = \arccos\!\left(\frac{\cot\theta}{\beta}\right),
\]
so
\[
  \cot\theta = \beta\cos(\phi - \phi_0).
\]
Expanding:
\[
  \cot\theta = \beta\cos\phi\cos\phi_0 + \beta\sin\phi\sin\phi_0.
\]
Multiplying through by $\sin\theta$ and using $x = a\sin\theta\cos\phi$, $y = a\sin\theta\sin\phi$, $z = a\cos\theta$:
\[
  \cos\theta = \beta\sin\theta\cos\phi\cos\phi_0 + \beta\sin\theta\sin\phi\sin\phi_0,
\]
\[
  a\cos\theta = \beta\cos\phi_0\,(a\sin\theta\cos\phi) + \beta\sin\phi_0\,(a\sin\theta\sin\phi),
\]
\[
  z = (\beta\cos\phi_0)\,x + (\beta\sin\phi_0)\,y.
\]
This can be written as
\[
  Ax + By - z = 0,
\]
where $A = \beta\cos\phi_0$ and $B = \beta\sin\phi_0$ are constants. This is the equation of a plane passing through the origin. Its intersection with the sphere $x^2 + y^2 + z^2 = a^2$ is a \textbf{great circle}.

\medskip
\noindent\textbf{Nature of the extremum:} A great circle arc that is less than a semicircle gives a minimum. However, if the arc is greater than a semicircle, it gives neither a minimum nor a maximum---it is a saddle point. For two non-antipodal points, the shorter great circle arc is a true minimum of the arc length; the longer arc is not a minimum (nearby paths are shorter) but still satisfies the Euler--Lagrange equation.

\end{document}
